% ============================================================
\chapter{Modul 5: Listen}
% ============================================================

\section{Lektion 15: Listen verstehen}

\begin{analogie}
Ein Laserscanner dreht sich 360 Grad und misst in jedem Grad den Abstand -- 360 Messwerte auf einmal. Diese Werte in 360 einzelnen Variablen zu speichern wäre absurd. Eine Liste speichert alle 360 Werte unter einem einzigen Namen.
\end{analogie}

\begin{lstlisting}
# Laserscanner: 8 Messwerte (vereinfacht)
abstaende = [0.8, 1.2, 2.4, 0.3, 1.8, 2.1, 0.9, 1.5]

# Zugriff per Index (beginnt bei 0)
print(abstaende[0])    # 0.8  -- erster Wert
print(abstaende[-1])   # 1.5  -- letzter Wert
print(abstaende[2:5])  # [2.4, 0.3, 1.8]  -- Slice

# Länge
print(f"Anzahl Messwerte: {len(abstaende)}")

# Suchen
print(f"Minimum: {min(abstaende)}")  # 0.3 -- nächstes Hindernis
print(f"Maximum: {max(abstaende)}")  # 2.4 -- freieste Richtung
\end{lstlisting}

\begin{merksatz}
Indices beginnen bei 0. Das letzte Element hat Index \texttt{len(liste) - 1} -- oder einfach \texttt{[-1]}. Slicing: \texttt{[start:stop]} gibt Elemente von start bis stop-1 zurück.
\end{merksatz}

\section{Lektion 16: Listen durchlaufen}

\begin{lstlisting}
abstaende = [0.8, 1.2, 0.15, 1.8, 0.3]

# Alle Hindernisse finden (< 0.25 m)
print("Hindernisse erkannt:")
for i, abstand in enumerate(abstaende):
    if abstand < 0.25:
        print(f"  Richtung {i * 45} Grad: {abstand} m")

# List Comprehension: elegante Filterung
hindernisse = [a for a in abstaende if a < 0.25]
sicher      = [a for a in abstaende if a >= 0.50]

print(f"Hindernis-Werte: {hindernisse}")
print(f"Freie Richtungen: {len(sicher)} von {len(abstaende)}")

# Durchschnitt berechnen
durchschnitt = sum(abstaende) / len(abstaende)
print(f"Mittlerer Abstand: {durchschnitt:.2f} m")
\end{lstlisting}

\section{Lektion 17: Listen verändern}

\begin{lstlisting}
log = []   # Leerer Log

# Sensormessungen hinzufügen
for i in range(5):
    import random
    messwert = round(random.uniform(0.2, 2.0), 2)
    log.append(messwert)
    print(f"Messung {i+1}: {messwert} m")

print(f"\nGesamter Log: {log}")

# Ältesten Wert entfernen (FIFO-Puffer -- typisch in Robotik)
alter_wert = log.pop(0)
print(f"Entfernter Wert: {alter_wert}")

# Sortieren: freieste Richtung finden
log_sortiert = sorted(log, reverse=True)
print(f"Beste Richtung: {log_sortiert[0]} m")

# Zusammenführen
neues_log = [1.2, 0.8, 1.5]
log.extend(neues_log)
\end{lstlisting}

\begin{praxis}
In der Robotik ist der \textbf{Ringpuffer} (FIFO: First In, First Out) fundamental: Der älteste Sensorwert fällt raus, der neueste kommt rein. Das glättet Ausreißer und ergibt stabilere Messungen. \texttt{append()} und \texttt{pop(0)} implementieren genau das.
\end{praxis}

\begin{uebung}[title=Projektaufgabe Modul 5: Laserscanner-Auswertung]
Simuliere einen 360-Grad-Laserscanner:
\begin{enumerate}
    \item Erzeuge eine Liste mit 36 Zufallswerten (0.1 bis 3.0 m): \texttt{[round(random.uniform(0.1, 3.0), 2) for \_ in range(36)]}
    \item Finde alle Richtungen mit Abstand $<$ 0.5 m (Hindernisse)
    \item Finde die Richtung mit dem größten Abstand (bester Weg)
    \item Berechne Minimum, Maximum, Durchschnitt
    \item Gib einen formatierten Scan-Bericht aus
\end{enumerate}
\end{uebung}
